\section*{Hoofdstuk \ref{ch:g12:continental-crust} --- De continentale korst}

\begin{solution}{exo:g12:continental-crust:1}
\omterm{def:g12:continental-crust:continental-crust}{Continentale korst}: 30 tot \qty{70}{km}, gemiddeld ongeveer \qty{41}{km};
granitische gesteenten, gneisen en andere metamorfe en plutonische gesteenten onder een
sedimentaire bedekking; gemiddeld ongeveer \qty{61}{\percent} silica; ongeveer
\qty{2830}{kg/m^3}; tot 4.03 miljard jaar (korrels 4.4). \omterm{def:g11:ocean-ridges:oceanic-lithosphere}{Oceanische korst}:
ongeveer \qty{7}{km}; basalt en \omterm{def:g11:ocean-ridges:gabbro}{gabbro}; ongeveer \qty{50}{\percent} silica; ongeveer
\qty{3000}{kg/m^3}; minder dan 200 miljoen jaar.
\end{solution}

\begin{solution}{exo:g12:continental-crust:2}
\qty{60.6}{\percent}: \omterm{def:g11:igneous-rocks:felsic-mafic}{intermediair}, zoals een \omterm{def:g11:igneous-rocks:rock-names}{dioriet}. \qty{66.6}{\percent}:
\omterm{def:g11:igneous-rocks:felsic-mafic}{felsisch}, zoals een graniet.
\end{solution}

\begin{solution}{exo:g12:continental-crust:3}
De \omterm{def:g11:ocean-ridges:oceanic-lithosphere}{oceanische korst}, ongeveer \qty{3000}{kg/m^3}, ligt op lithosferische mantel,
en de oude, koude \omterm{def:g11:ocean-ridges:oceanic-lithosphere}{oceanische lithosfeer} als geheel is dicht genoeg om te zinken.
De \omterm{def:g12:continental-crust:continental-crust}{continentale korst}, ongeveer \qty{2830}{kg/m^3}, veel lichter dan de mantel
(ongeveer \qty{3380}{kg/m^3}), heeft te veel opwaartse kracht: een continent dat bij een trog aankomt,
stopt de \omterm{def:g11:subduction-zones:subduction}{subductie}.
\end{solution}

\begin{solution}{exo:g12:continental-crust:4}
(1) Een oude, stabiele kern van een continent, die niet meer vervormt. (2) Meer
dan een miljard jaar. (3) Metamorfe en plutonische gesteenten, de diepe \omterm{def:g12:continental-crust:root}{wortels} van
vroegere gebergten.
\end{solution}

\begin{solution}{exo:g12:continental-crust:5}
$4030/190 \approx 21$ keer. De oceaanbodem wordt bij de ruggen gemaakt en
bij \omterm{def:g11:subduction-zones:subduction}{subductiezones} vernietigd: niets ervan overleeft veel langer dan 200 miljoen
jaar.
\end{solution}

\begin{solution}{exo:g12:continental-crust:6}
(1) $0.306 \times 210.4 \approx \num{64.4}\text{ miljoen }\unit{km^2}$.
(2) $210.4 - 64.4 = \num{146.0}\text{ miljoen }\unit{km^2}$.
(3) $4\pi \times 6371^2 \approx \qty{5.10e8}{km^2}$; $210.4/510.1 \approx
0.41$: \qty{41}{\percent}.
\end{solution}

\begin{solution}{exo:g12:continental-crust:7}
(1) $41100 \times 2830 \approx \qty{1.16e8}{kg}$. (2) $7100 \times 3000 =
\qty{2.13e7}{kg}$. (3) Ongeveer 5.5 keer.
\end{solution}

\begin{solution}{exo:g12:continental-crust:8}
\omterm{def:g12:continental-crust:root}{Wortel}: $2.5 \times 2830/550 \approx \qty{12.9}{km}$. De korst is dikker met
$2.5 \times 3380/550 \approx \qty{15.4}{km}$: in totaal ongeveer \qty{50}{km},
binnen de 50 tot \qty{70}{km} die onder gebergtegordels is gevonden.
\end{solution}

\begin{solution}{exo:g12:continental-crust:9}
(1) Ongeveer $0.103 + 0.100 + 0.102 \approx 0.30$: \qty{30}{\percent}.
(2) De jongste band, 0 tot 10 miljoen jaar (ongeveer \qty{10}{\percent}),
bijna gelijk met 20 tot 30. (3) Oude bodem is bij \omterm{def:g11:subduction-zones:subduction}{subductiezones} vernietigd;
hoe ouder de bodem, hoe meer ervan verdwenen is.
\end{solution}

\begin{solution}{exo:g12:continental-crust:10}
$120 = 0.694\,h - 0.306 \times 100$, dus $h = 150.6/0.694 \approx
\qty{217}{m}$. Wie alleen het land telt, laat de lage, ondergelopen randen van
de continenten weg, en het gemiddelde valt veel hoger uit.
\end{solution}

\begin{solution}{exo:g12:continental-crust:11}
Het eiland: een \omterm{def:g11:ocean-ridges:gabbro}{gabbro} onder lava's en een \omterm{def:g11:seismic-structure:moho}{Moho} op \qty{8}{km}, \omterm{def:g11:ocean-ridges:oceanic-lithosphere}{oceanische korst}
(de hoogte van het eiland verandert daar niets aan). De vlakte: gneis en een
\omterm{def:g11:seismic-structure:moho}{Moho} op \qty{38}{km}, \omterm{def:g12:continental-crust:continental-crust}{continentale korst}.
\end{solution}

\begin{solution}{exo:g12:continental-crust:12}
(1) Exact: $\tfrac{4}{3}\pi(6371^3 - 6329.9^3) \approx
\qty{2.0829e10}{km^3}$; benaderd: $4\pi \times 6371^2 \times 41.1
\approx \qty{2.0964e10}{km^3}$. (2) De benadering is ongeveer
\qty{0.65}{\percent} te groot, dicht bij $H/R \approx 0.0065$: voor een dunne laag is de
fout verwaarloosbaar.
\end{solution}

\begin{solution}{exo:g12:continental-crust:13}
(1) $60.6 = (66.6 + 2x)/3$, dus $x = (181.8 - 66.6)/2 = \qty{57.6}{\percent}$.
(2) \omterm{def:g11:igneous-rocks:felsic-mafic}{Intermediair}. (3) De diepe \omterm{def:g12:continental-crust:continental-crust}{continentale korst} is minder rijk aan silica dan
de bovenkorst: meer mafische mineralen, rijker aan ijzer, magnesium en calcium.
\end{solution}

\begin{solution}{exo:g12:continental-crust:14}
(1) $3 \times 10^8/3 = 10^8$ jaar, 100 miljoen jaar. (2) Dezelfde orde als
de ouderdommen van de oceaanbodem: drie kwart jonger dan 100 miljoen jaar,
niets ouder dan ongeveer 190. (3) Ze moet even snel worden vernietigd als ze wordt gemaakt:
ze zinkt bij \omterm{def:g11:subduction-zones:subduction}{subductiezones} terug in de mantel.
\end{solution}

\begin{solution}{exo:g12:continental-crust:15}
(1) $8.6 \times 10^9 / 4.0 \times 10^9 \approx \qty{2.2}{km^3}$ per jaar.
(2) Een constant tempo is maar één mogelijkheid: de korst kan vroeg snel zijn gegroeid
en sindsdien gerecycled, of laat zijn gegroeid. De oudste gesteenten en korrels bewijzen
dat er vroeg al wat korst bestond; ze kunnen niet zeggen hoeveel, omdat oude korst
kan worden geërodeerd of teruggevoerd in de mantel en weinig sporen kan achterlaten.
\end{solution}

\begin{solution}{pb:g12:continental-crust:1}
\textbf{1.} $4\pi \times 6371^2 \approx \qty{5.10e8}{km^2}$.

\textbf{2.} $210.4/510.1 \approx 0.41$: \qty{41}{\percent}.

\textbf{3.} $0.306 \times 210.4 \approx \num{64.4}\text{ miljoen }\unit{km^2}$.

\textbf{4.} $510.1 - 210.4 \approx \num{299.7}\text{ miljoen }\unit{km^2}$,
\qty{59}{\percent}.

\textbf{5.} De randen van de continenten liggen onder ondiepe zeeën, op de
continentale platten.

\textbf{6.} De dikte, \qty{41}{km}, is heel klein vergeleken met de
straal, \qty{6371}{km}: de kromming verandert het volume met minder dan één
procent.

\textbf{7.} $210.4 \times 10^6 \times 41.1 \approx \qty{8.65e9}{km^3}$.

\textbf{8.} \qty{8.65e18}{m^3}.

\textbf{9.} $210.4 \times 10^6 \times 36 \approx \qty{7.57e9}{km^3}$.

\textbf{10.} $(8.65 - 7.57)/7.57 \approx 0.14$: \qty{14}{\percent}.

\textbf{11.} $8.65 \times 10^{18} \times 2830 \approx \qty{2.45e22}{kg}$.

\textbf{12.} $41100 \times 2830 \approx \qty{1.16e8}{kg}$.

\textbf{13.} $299.7 \times 10^6 \times 7.1 \approx \qty{2.13e9}{km^3}$, dat
is \qty{2.13e18}{m^3}; maal 3000: ongeveer \qty{6.4e21}{kg}.

\textbf{14.} De \omterm{def:g12:continental-crust:continental-crust}{continentale korst}, ongeveer 3.8 keer zwaarder: ze is bijna zes
keer dikker, en dat weegt zwaarder dan haar kleinere oppervlakte en dichtheid.

\textbf{15.} Ze bestaat uit granitische gesteenten, rijk aan kwarts en veldspaten, de
lichtste gesteentevormende mineralen; de \omterm{def:g11:ocean-ridges:oceanic-lithosphere}{oceanische korst} bestaat uit basalt en
\omterm{def:g11:ocean-ridges:gabbro}{gabbro}, rijk aan pyroxeen en olivijn.

\textbf{16.} $2.45 \times 10^{22}/5.97 \times 10^{24} \approx 0.0041$.

\textbf{17.} $3 \times 0.41 \times (41.1/6371) \times (2830/5513) \approx
0.0041$: hetzelfde.

\textbf{18.} $6.4 \times 10^{21}/5.97 \times 10^{24} \approx 0.0011$: ongeveer een
tiende procent.

\textbf{19.} Ongeveer \qty{0.4}{\percent}.

\textbf{20.} De \omterm{def:g12:continental-crust:continental-crust}{continentale korst} is ongeveer \textbf{\qty{0.4}{\percent} van de
massa van de Aarde}.
\end{solution}
